\chapter[Sermon 22. He Draws Them Also Hereby unto Repentance.]{He Draws Them Also Hereby unto Repentance.}
\label{sermon:022}
\markright{Sermon 22. He Draws Them Also Hereby unto Repentance.}

Behold, I stand at the door and knock. If anyone hears my voice and opens the door, I will come in to him, and will share a meal with him, and he with me. To the one who overcomes, I will grant the privilege of sitting with me on my throne, just as I overcame and have sat with my Father on his throne. Let the one who has ears hear what the Spirit says to the congregations.

\index{Repentance}Hereby also the Lord allures the Laodiceans to repentance, showing that every time is suitable for conversion, and that God is ever ready to receive sinners, and constantly urges them to repent and live. And this matter he explains in an allegorical and beautiful speech, taken from the fifth chapter of the book of Song of Songs. For he imagines the Lord standing at the door and knocking, indeed promising to those who open it the greatest intimacy and unspeakable joys.

First, therefore, is declared the benevolence of God toward sinners, and his most ready will always to receive them, yes, and his infinite study to move men to repentance, that they might live. For the Lord stands at the door, and knocks. The word “stands” signifies that God is always prepared, always watches over our salvation. For he sits not idly, nor lies on one side like a sluggard: He stands busily to his work. And “I stand,” he says, not “I stood” or “shall stand,” but “I stand” evermore ready, evermore loving and gentle. What does he do? He knocks, and that indeed at the door, desiring to be let in. For just as he who knocks at a door seriously covets to be let in, so God desires earnestly to be received by us. And God uses sundry kinds of knocking. For he warns and excites with his word by the Prophets, again by signs and wonders, and also by sundry chances and movements [?]. These things may be seen in the city of Jerusalem. He sends to them his Prophets and Apostles. He shows divers wonders. He brings on the sorrowful chances that they might admonish them—such as are reported in Luke 13, of the Galileans and of those whom the tower of Siloam had overwhelmed. We may see the like at this day, how the Lord knocks. Therefore he said truly, “Jerusalem, Jerusalem,” and so forth, Matthew 23. These are doubtless the parts and doings of God who will not that a sinner should die, but rather convert and live.

\index{Hebrews, Epistle to the}Then we must see what is required of us, namely that we should hear the knocking and noise of the knocker, and also open and receive those who desire to come in. Here are refuted those who speak of humankind as though it were a block, and an image of some unknown design, saying: It is neither in the runner nor in the willer, and so on. Some entirely abstain from doing good, saying, “If I am chosen, that is enough.” But Scripture requires everywhere hearing and obedience. We know that people are only saved, and that in Christ: in Christ to be those who believe; that faith is of hearing, hearing by the word of God. Therefore, the prophet says, “This day if you hear his voice.” This is recited by the Apostle in Hebrews 4. The Apostle also, in 2 Timothy 2, says, “In a great house there are not only vessels of gold, but of earth also.” If anyone purges himself [?], and so on. And therefore the Lord says, “I knock.” It shall be your part, not to despise him who knocks, but to open to him. And he mentions indeed two things: to hear, which is required both in John 8 and John 9 of the children of God and the true sheep; and to open, that is to receive the Lord, or believe, to obey, and to shape oneself after the will of God, and to do penance. Nevertheless, we must be on our guard here, so that we do not think that humankind has power of itself to receive the Lord. The Lord illuminates his elect, and by him we can do all things; without him we can do nothing. Other passages must be consulted with this, such as John 15:2, 1 Corinthians 3, and Philippians 2. Therefore, those who open, open by the grace of God. Those who do not open are wrapped in their sins, not through their own fault, but through their own fault they do not open, and not through any fault of God.

\index{Turks}Let us hear moreover what the Lord promises to those who open [their hearts], that is to say, to such as receive Christ with true faith. The Lord promises them two things chiefly. First, “I will go into him,” he says. Scripture signifies that Christ dwells everywhere through faith in the hearts of the faithful, and with a most tight bond to be joined unto them. “He who eats my flesh and drinks my blood abides in me, and I in him.” These things are spoken of the Lord in the sixth chapter of John. And in the fourteenth chapter, he says, “He who loves me will keep my word, and my Father and I will come unto him and will make abode with him.” Paul says that he no longer lives, but that Christ lives in him. He affirms that Christ through faith dwells in the hearts of the faithful. And so the Lord enters the hearts of those who let him in. Not the least part of felicity consists in this conjunction. For to be united with God is blessedness, which begins here and is made perfect in another life. And therefore, in the second place, the Lord says, “And I will sup with him, and he with me.” Whereby he notes not only again a most dear friendship and familiarity (for the table is consecrated to amity) but rather the fruition of eternal glory. For by the supper are signified the celestial joys, greatest and unspeakable, which the godly receive immediately after death, but more fully in the end of times, when the bodies shall arise again. Therefore it is not applied to a dinner, but to a supper, as it is also in the fourteenth chapter of Luke. If we receive Christ, we shall have him dwelling with us continually while we live in this world. And in the world to come we shall have the full fruition of all the celestial joys. These things are certain and true. For otherwise in the life to come there shall be no riotous banquets, such as the Turks do imagine.

He adds another general promise, by which he exhorts and encourages the study of godly religion and repentance. For to him who overcomes is promised the kingdom of heaven. And he says to him who overcomes—of which I have spoken in other epistles—not to him who flees or is a coward, etc. He presents the example of conquering Christ. For we must overcome, as he has overcome. He indeed overcame most perfectly; we, with our limited strength, fight and overcome. And truly, the true victory in us is the lively virtue of Christ: that is to say, by him they overcome, whoever overcomes. And just as he, having overcome death and vanquished the world and the devil, ascended into heaven and sat on the right hand of the Father, so he promises us, upon overcoming, that he will give us the seat of his Father: not that we, sitting on the right hand of God, should judge over all flesh, being made Christ’s, but that being made partakers of everlasting glory and delivered from all judgment, we may appear in glory when he shall come to judge the living and the dead. We read of a similar promise made to the disciples in Matthew 19 and Luke 22. And this glory will assuredly come to us, just as Christ himself did truly ascend into heaven and sit in celestial glory.

And here we must note a special thing, that Christ gives here that which in the twentieth chapter of Matthew he denies he can give to James and John, namely, to sit in celestial glory. Therefore this passage explains that. For Christ, after declaring his deity, gives what he later, in his humanity, denies he can give. This passage then proves that Christ is very God, giver of eternal life, and so forth.

He then, in his characteristic way, adds an affirmation, by which he applies this letter to all congregations, and asserts that it is inspired by the Spirit of Christ. Concerning this, we have spoken previously.

And we have treated thus far of the second part of this work, wherein are declared the most excellent points of our religion, namely, who and of what sort is Christ, sitting in the glory of the Father, how he is present in his church, and governs it as king and priest, by his Spirit, by his word, and by the Sacraments. Also, what and of what sort is the church of Christ? What is the true and right doctrine of the church? What opinions are wicked? What is to be done with erroneous doctrines and seducers? How may the church, having fallen and been afflicted, be restored? What is true repentance, and what are the duties of the godly, and many other things of like sort. To God the Father be praise, thanksgiving, and glory, through Jesus Christ our Lord.
